
\documentclass[11pt,a4paper]{article}

\usepackage[a4paper,margin=2cm]{geometry}
\usepackage{fontspec}
\usepackage{polyglossia}
\setdefaultlanguage{vietnamese}
\setmainfont{TeX Gyre Heros}

\usepackage{xcolor}
\usepackage{array}
\usepackage{tabularx}
\usepackage{booktabs}
\usepackage{colortbl}
\usepackage{enumitem}
\usepackage{fancyhdr}
\usepackage{ragged2e}
\usepackage{setspace}
\usepackage{titlesec}

\definecolor{titleblue}{RGB}{35,70,175}
\definecolor{tableblue}{RGB}{235,243,253}

\setlength{\parindent}{0pt}
\setlength{\parskip}{4pt}
\renewcommand{\arraystretch}{1.35}
\setlength{\tabcolsep}{5pt}

\pagestyle{fancy}
\fancyhf{}
\fancyfoot[C]{\thepage}
\renewcommand{\headrulewidth}{0pt}
\renewcommand{\footrulewidth}{0pt}

\titleformat{\section}
{\color{titleblue}\Large\bfseries}
{\thesection.}{0.4em}{}

\titlespacing*{\section}{0pt}{12pt}{16pt}

\setlist[enumerate]{
    leftmargin=1.1cm,
    itemsep=6pt,
    topsep=5pt
}

\begin{document}

\section*{1. Product Vision (Định hướng sản phẩm)}

\vspace{8pt}

\begin{center}
\fbox{
\begin{minipage}{0.93\textwidth}
\centering
\large\itshape
" Hệ Thống Đặt Và Thanh Toán Đồ Uống giúp khách hàng
đặt và thanh toán đồ uống không tiền mặt chỉ trong vài
chạm, mang lại trải nghiệm nhanh chóng, tiện lợi và
hiện đại trong không gian quán. "
\end{minipage}
}
\end{center}

\vspace{55pt}

\section*{2. Bảng Stakeholders (Các bên liên quan)}

\vspace{12pt}

\small
\begin{tabularx}{\textwidth}{
@{}>{\RaggedRight\arraybackslash\bfseries}p{2.3cm}
>{\RaggedRight\arraybackslash}p{2.6cm}
>{\RaggedRight\arraybackslash}X
>{\RaggedRight\arraybackslash}X@{}
}
\toprule
\rowcolor{tableblue}
Stakeholder & Vai trò & Mục tiêu / Kỳ vọng &
Mối quan tâm / Thách thức \\
\midrule

Khách hàng &
Người dùng cuối &
Xem menu nhanh, tùy chọn size/topping dễ dàng,
đặt hàng và thanh toán không tiền mặt không lo xếp hàng. &
Trải nghiệm mượt mà, không giật lag giờ cao điểm,
tính chính xác của đơn hàng. \\

\midrule

Quản lý / PO &
Quản lý vận hành &
Quản lý món ăn, tồn kho, theo dõi trạng thái đơn hàng
thời gian thực và tối ưu quy trình phục vụ. &
Giao diện quản trị trực quan, báo cáo doanh thu chính xác,
giảm thiểu sai sót đơn. \\

\midrule

Giảng viên &
Nghiệm thu &
Kiểm tra tính đầy đủ của các yêu cầu chức năng,
kiến trúc hệ thống và khả năng vận hành thực tế. &
Tuân thủ quy trình phát triển phần mềm, chuẩn hóa code,
tính bảo mật và tài liệu kỹ thuật. \\

\midrule

Dev Team &
Phát triển hệ thống &
Phát triển tính năng đúng tiến độ, mã nguồn sạch,
kiến trúc dễ nâng cấp và duy trì. &
Khả năng mở rộng, quản lý API, tính đồng bộ giữa
Frontend và Backend. \\

\midrule

Gateway Mock &
Trung gian thanh toán &
Xử lý giao dịch thanh toán trực tuyến an toàn,
chính xác và phản hồi nhanh chóng. &
Tính bảo mật thông tin thanh toán, xử lý phản hồi
giao dịch (Idempotency). \\

\bottomrule
\end{tabularx}
\normalsize

\newpage

\section*{3. Yêu cầu chức năng (Functional Requirements - FR)}

\vspace{12pt}

\small
\begin{tabularx}{\textwidth}{
@{}p{0.8cm}
>{\RaggedRight\arraybackslash\bfseries}p{2.1cm}
>{\RaggedRight\arraybackslash}X
>{\RaggedRight\arraybackslash}p{2.0cm}@{}
}
\toprule
\rowcolor{tableblue}
ID & Đối tượng & Yêu cầu chức năng (FR) &
Tương ứng nhánh BFD \\
\midrule

FR-01 &
Tất cả người dùng &
Hệ thống yêu cầu đăng nhập, đăng xuất và cho phép
cập nhật thông tin hồ sơ cá nhân. &
1. Quản lý tài khoản \\

\midrule

FR-02 &
Quản lý (Admin) &
Hệ thống cho phép quản trị viên thêm, sửa, xóa
tài khoản nhân viên và phân quyền truy cập. &
1. Quản lý tài khoản \\

\midrule

FR-03 &
Quản lý (Admin) &
Hệ thống cho phép thêm, sửa, xóa, phân loại danh mục
và cập nhật trạng thái món ăn (còn/hết). &
2. Quản lý thực đơn \\

\midrule

FR-04 &
Quản lý (Admin) &
Hệ thống hỗ trợ tạo phiếu nhập kho, kiểm kê và
cập nhật số lượng tồn kho nguyên liệu. &
3. Quản lý nhập kho \\

\midrule

FR-05 &
Khách / Phục vụ &
Hệ thống cho phép xem chi tiết thực đơn, tìm kiếm
món ăn và xem tư vấn món bán chạy. &
4. Đặt món \\

\midrule

FR-06 &
Khách / Phục vụ &
Hệ thống cho phép thêm món vào giỏ hàng, tùy chọn
hình thức (tại bàn/mang đi), thêm ghi chú và
xác nhận gửi order. &
4. Đặt món \\

\bottomrule
\end{tabularx}
\normalsize

\newpage

\section*{4. Yêu cầu phi chức năng (Non-Functional Requirements - NFR)}

\vspace{12pt}

\small
\begin{tabularx}{\textwidth}{
@{}p{0.8cm}
>{\RaggedRight\arraybackslash\bfseries}p{2.2cm}
>{\RaggedRight\arraybackslash}X@{}
}
\toprule
\rowcolor{tableblue}
ID & Loại & Non-functional Requirement \\
\midrule

NFR-01 &
Performance &
Ít nhất 95\% yêu cầu xem danh sách món ăn phải được hệ thống
phản hồi trong vòng 2 giây trong điều kiện tải bình thường
(Thời gian phản hồi API dưới 500ms). \\

\midrule

NFR-02 &
Performance &
Hệ thống phải có khả năng xử lý ít nhất 100 yêu cầu đặt món
đồng thời mà không làm mất, tạo trùng hoặc sai thông tin đơn hàng. \\

\midrule

NFR-03 &
Security &
Hệ thống phải yêu cầu khách hàng đăng nhập trước khi thực hiện
các thao tác đặt món và thanh toán. \\

\midrule

NFR-04 &
Security &
Hệ thống phải phân quyền theo vai trò (RBAC), đảm bảo khách hàng
không thể truy cập các chức năng quản lý dành cho nhân viên
và quản trị viên. \\

\midrule

NFR-05 &
Security &
Thông tin tài khoản và mật khẩu của người dùng phải được bảo vệ
bằng mã hóa băm (bcrypt) và token JWT, không cho phép người dùng
khác truy cập trái phép. \\

\midrule

NFR-06 &
Usability &
Giao diện đặt món phải rõ ràng, dễ sử dụng và cho phép khách hàng
hoàn thành quy trình chọn món $\rightarrow$ thêm vào giỏ hàng
$\rightarrow$ xác nhận đơn hàng một cách thuận tiện. \\

\midrule

NFR-07 &
Reliability &
Hệ thống phải đảm bảo thông tin về món ăn, số lượng, giá tiền
và đơn hàng được lưu trữ chính xác, không bị mất hoặc sai lệch
trong quá trình xử lý. \\

\midrule

NFR-08 &
Availability &
Hệ thống phải duy trì hoạt động ổn định trong thời gian phục vụ
và hạn chế tối đa tình trạng gián đoạn khi khách hàng thực hiện
đặt món. \\

\midrule

NFR-09 &
Scalability &
Hệ thống phải cho phép thêm mới cửa hàng, món ăn, nhân viên
và tài khoản quản lý mà không ảnh hưởng đến các chức năng
đang hoạt động. \\

\midrule

NFR-10 &
Maintainability &
Hệ thống phải được thiết kế và tổ chức mã nguồn rõ ràng,
thuận tiện cho việc bảo trì và mở rộng tính năng về sau. \\

\bottomrule
\end{tabularx}

\normalsize

\vspace{25pt}

\section*{5. User Story \& Acceptance Criteria tiêu biểu}

\vspace{8pt}

\textbf{\large US-05: Thanh toán không tiền mặt}

\vspace{12pt}

\hspace{1cm}\textbf{User Story:}

\hspace{1cm}
"Là một Khách hàng, tôi muốn thực hiện thanh toán đơn hàng
bằng phương thức không tiền mặt (chuyển khoản / cổng thanh toán)
để tôi có thể hoàn tất việc đặt đồ uống nhanh chóng mà không
cần đợi trả lại tiền mặt."

\vspace{15pt}

\textbf{Acceptance Criteria (Tiêu chí chấp nhận):}

\begin{enumerate}

\item \textbf{Given} Khách hàng đã chọn xong các món trong giỏ hàng
và tiến hành sang bước thanh toán.

\item \textbf{When} Khách hàng chọn phương thức
"Thanh toán qua Ví điện tử / Mã QR" và nhấn
"Xác nhận thanh toán".

\item \textbf{Then} Hệ thống tạo mã giao dịch duy nhất,
hiển thị thông tin số tiền và mã QR thanh toán tương ứng
trong vòng dưới 500ms.

\end{enumerate}

\newpage

\vspace*{0pt}

\begin{enumerate}[start=4]

\item \textbf{And} Khi giao dịch thanh toán thành công,
trạng thái đơn hàng lập tức được cập nhật thành
\texttt{PAID} và gửi thông báo xác nhận đơn hàng
cho khách hàng.

\item \textbf{And} Trường hợp thanh toán thất bại hoặc quá
thời gian chờ (timeout), hệ thống thông báo lỗi rõ ràng
và giữ nguyên giỏ hàng để khách hàng thử lại mà không
bị trừ tiền hoặc tạo trùng đơn.

\end{enumerate}

\end{document}
